% 对应单元：QM-C01-S01--S03
% 来源：Shankar 第 1 章第 1.1--1.3 节及同书练习
% 人工核验：用户作答后完成订正与收口检查

\section{第 1.1 节教材习题与订正}
\label{exercise-set:QM-C01-S01}

\exercisemeta
  {QM-C01-S01-EX}
  {2026-08-17}
  {Shankar，第 1 章 \S 1.1，练习 1.1.1--1.1.5}
  {全部完成并逐题核对}
  {已确认（文字与手写解答）}

\exercisequestion{QM-C01-S01-E01}{教材练习 1.1.1：由公理推出四个基本性质}

\textbf{对应知识点：}\relatedtopic{QM-C01-S01-T01}{向量空间、标量域与公理}。

\textbf{题意概述。}证明零向量唯一、$0\ket{V}=\ket{0}$、$(-1)\ket{V}=-\ket{V}$，以及加法逆元唯一。教材位置：印刷页 p.~2，PDF p.~16。

\begin{checkedsolution}
若 $\ket{0}$ 与 $\ket{0'}$ 都是零向量，则
\[
  \ket{0'}=\ket{0'}+\ket{0}
  =\ket{0}+\ket{0'}=\ket{0},
\]
故零向量唯一。

由
\[
  0\ket{V}=(0+0)\ket{V}=0\ket{V}+0\ket{V}
\]
在等式两边加上 $0\ket{V}$ 的加法逆元，得到 $0\ket{V}=\ket{0}$。

又有
\[
  \ket{V}+(-1)\ket{V}
  =1\ket{V}+(-1)\ket{V}
  =(1-1)\ket{V}
  =0\ket{V}=\ket{0}.
\]
因此 $(-1)\ket{V}$ 是 $\ket{V}$ 的加法逆元，即 $(-1)\ket{V}=-\ket{V}$。

最后，若 $\ket{W}$ 也满足 $\ket{V}+\ket{W}=\ket{0}$，则
\[
  \ket{W}=\ket{W}+\ket{0}
  =\ket{W}+(\ket{V}-\ket{V})
  =(\ket{W}+\ket{V})-\ket{V}
  =\ket{0}-\ket{V}=-\ket{V}.
\]
故加法逆元唯一。
\end{checkedsolution}

\textbf{检查点：}零向量 $\ket{0}$ 与标量零 $0$ 是不同类型的对象；证明中使用 $0\ket{V}=\ket{0}$ 时不可混写。

\exercisequestion{QM-C01-S01-E02}{教材练习 1.1.2：实三元组与非封闭子集}

\textbf{对应知识点：}\relatedtopic{QM-C01-S01-T02}{典型例子、封闭性与齐次约束}。

\textbf{题意概述。}实三元组按分量相加和数乘。写出零向量与 $(a,b,c)$ 的加法逆元，并说明形如 $(a,b,1)$ 的三元组不构成向量空间。教材位置：印刷页 p.~3，PDF p.~17。

\begin{checkedsolution}
零向量为
\[
  (0,0,0),
\]
而 $(a,b,c)$ 的加法逆元为
\[
  (-a,-b,-c).
\]
集合 $S=\{(a,b,1):a,b\in\real\}$ 不包含零向量。若 $\bm{u},\bm{v}\in S$，则 $\bm{u}+\bm{v}$ 的第三分量为 $2$，一般不在 $S$；若 $\lambda\in\real$，则 $\lambda\bm{u}$ 的第三分量为 $\lambda$，一般也不在 $S$。此外，$-\bm{u}$ 的第三分量为 $-1$。因此该集合不构成向量空间。
\end{checkedsolution}

\textbf{检查点：}任意一个必要条件失败就足以否定向量空间；列出零向量和封闭性失败能更完整地定位问题。

\exercisequestion{QM-C01-S01-E03}{教材练习 1.1.3：带边界条件的函数集合}

\textbf{对应知识点：}\relatedtopic{QM-C01-S01-T02}{典型例子、封闭性与齐次约束}。

\textbf{题意概述。}判断满足端点为零、周期边界条件以及 $f(0)=4$ 的函数集合是否构成向量空间。教材位置：印刷页 p.~4，PDF p.~18。

\begin{checkedsolution}
若 $f(0)=f(L)=0$ 且 $g(0)=g(L)=0$，则 $f+g$ 与 $af$ 仍在两个端点取零；零函数和 $-f$ 也满足条件。因此该集合构成向量空间。

若 $f(0)=f(L)$ 且 $g(0)=g(L)$，则
\[
  (f+g)(0)=(f+g)(L),\qquad
  (af)(0)=(af)(L),
\]
零函数和加法逆函数也满足条件，因此周期函数集合构成向量空间。

对于 $f(0)=4$，零函数不在集合中；两个集合元素相加后在 $0$ 处取值为 $8$；一般数乘后在 $0$ 处取值为 $4a$；取负后为 $-4$。因此该集合不构成向量空间。
\end{checkedsolution}

\textbf{检查点：}严格说法是“零函数不属于该集合”，而不是“零函数不存在”。

\exercisequestion{QM-C01-S01-E04}{教材练习 1.1.4：三个矩阵的线性相关性}

\textbf{对应知识点：}\relatedtopic{QM-C01-S01-T03}{线性组合、线性相关与线性无关}。

\textbf{题意概述。}判断以下三个实矩阵是否线性无关：
\[
  \ket{1}=\begin{pmatrix}0&1\\0&0\end{pmatrix},\qquad
  \ket{2}=\begin{pmatrix}1&1\\0&1\end{pmatrix},\qquad
  \ket{3}=\begin{pmatrix}-2&-1\\0&-2\end{pmatrix}.
\]
教材位置：印刷页 p.~5，PDF p.~19。

\begin{checkedsolution}
直接计算可见
\[
  \ket{1}-2\ket{2}=\ket{3},
\]
所以
\[
  \ket{1}-2\ket{2}-\ket{3}=\ket{0}.
\]
系数 $(1,-2,-1)$ 不全为零，因此三个矩阵线性相关。
\end{checkedsolution}

\textbf{检查点：}证明线性相关不必求出所有线性关系；找到一条非平凡关系已经足够。

\exercisequestion{QM-C01-S01-E05}{教材练习 1.1.5：两组三维行向量}

\textbf{对应知识点：}\relatedtopic{QM-C01-S01-T03}{线性组合、线性相关与线性无关}。

\textbf{题意概述。}证明 $(1,1,0),(1,0,1),(3,2,1)$ 线性相关，并证明 $(1,1,0),(1,0,1),(0,1,1)$ 线性无关。教材位置：印刷页 p.~5，PDF p.~19。

\begin{checkedsolution}
对第一组，不难验证
\[
  2(1,1,0)+(1,0,1)=(3,2,1),
\]
故存在非平凡线性关系，第一组线性相关。

对第二组，令
\[
  a(1,1,0)+b(1,0,1)+c(0,1,1)=(0,0,0),
\]
比较分量得到
\[
  a+b=0,\qquad a+c=0,\qquad b+c=0.
\]
由这些等式可得 $a=b=c$，再代入 $a+b=0$ 得 $2a=0$。由于标量域为 $\real$，所以 $a=b=c=0$，第二组线性无关。
\end{checkedsolution}

\textbf{检查点：}相关性只需一个非平凡解；无关性要求证明齐次方程只有平凡解。

\exercisequestion{QM-C01-S01-E06}{自拟检查题：分量运算与换基}

\textbf{对应知识点：}\relatedtopic{QM-C01-S01-T05}{坐标唯一性与换基}。

\textbf{来源说明。}本题由学习对话生成，不是教材习题。它用于检查坐标运算、换基和“向量不随坐标改变”的区别。

在基 $\{\ket{1},\ket{2}\}$ 下，设
\[
  \ket{V}=3\ket{1}-2\ket{2},\qquad
  \ket{W}=-\ket{1}+4\ket{2}.
\]
求 $2\ket{V}+\ket{W}$ 的坐标。再令
\[
  \ket{1'}=\ket{1}+\ket{2},\qquad
  \ket{2'}=\ket{1}-\ket{2},
\]
求 $\ket{V}$ 在新基下的坐标。

\begin{checkedsolution}
第一问有
\[
  2\ket{V}+\ket{W}=5\ket{1},
\]
所以在原基下的坐标为 $(5,0)$。

设 $\ket{V}=a\ket{1'}+b\ket{2'}$。展开后
\[
  \ket{V}=(a+b)\ket{1}+(a-b)\ket{2},
\]
故
\[
  a+b=3,\qquad a-b=-2.
\]
解得
\[
  a=\frac12,\qquad b=\frac52.
\]
新坐标为 $(1/2,5/2)$；改变的只是坐标，抽象向量 $\ket{V}$ 本身没有改变。
\end{checkedsolution}

\textbf{检查点：}换基后将新坐标代回新基，必须能够重建原来的向量。

\section{第 1.2 节相关教材习题与订正}
\label{exercise-set:QM-C01-S02}

\exercisemeta
  {QM-C01-S02-EX}
  {2026-08-18}
  {Shankar，第 1 章 \S 1.2 Theorem 3 与 Exercise 1.3.1}
  {完成思路、关键投影计算和正交归一检查}
  {已确认（手写解答与收口问答）}

\exercisequestion{QM-C01-S02-E01}{教材练习 1.3.1：二维 Gram--Schmidt 正交化}

\textbf{对应知识点：}\relatedtopic{QM-C01-S02-T05}{Gram--Schmidt 正交化}。

\textbf{来源说明。}第 1.2 节本身没有编号习题；本题位于紧邻的同书后文，直接练习第 1.2 节的 Gram--Schmidt 定理。教材位置：印刷页 p.~15，PDF p.~29。

\textbf{题意概述。}对
\[
  \bm A=\begin{pmatrix}3\\4\end{pmatrix},
  \qquad
  \bm B=\begin{pmatrix}2\\-6\end{pmatrix},
\]
分别以 $\bm A$ 和 $\bm B$ 为第一个方向，构造两组正交归一基。

\begin{checkedsolution}
先以 $\bm A$ 为第一个方向：
\[
  \ket{e_1}=\frac{1}{5}\begin{pmatrix}3\\4\end{pmatrix},
  \qquad
  \braket{e_1}{B}=-\frac{18}{5}.
\]
减去 $\bm B$ 在 $\ket{e_1}$ 方向的投影：
\begin{align*}
  \ket{B_\perp}
  &=\ket{B}-\ket{e_1}\braket{e_1}{B}\\
  &=\begin{pmatrix}104/25\\-78/25\end{pmatrix}
   =\frac{26}{25}\begin{pmatrix}4\\-3\end{pmatrix}.
\end{align*}
因此第一组正交归一基可取为
\[
  \boxed{
  \ket{e_1}=\frac{1}{5}\begin{pmatrix}3\\4\end{pmatrix},
  \qquad
  \ket{e_2}=\frac{1}{5}\begin{pmatrix}4\\-3\end{pmatrix}}
\]

交换顺序，以 $\bm B$ 为第一个方向：
\[
  \ket{f_1}=\frac{1}{\sqrt{10}}\begin{pmatrix}1\\-3\end{pmatrix},
  \qquad
  \braket{f_1}{A}=-\frac{9}{\sqrt{10}}.
\]
相应的剩余向量为
\[
  \ket{A_\perp}
  =\ket{A}-\ket{f_1}\braket{f_1}{A}
  =\begin{pmatrix}39/10\\13/10\end{pmatrix}
  =\frac{13}{10}\begin{pmatrix}3\\1\end{pmatrix}.
\]
故第二组可取为
\[
  \boxed{
  \ket{f_1}=\frac{1}{\sqrt{10}}\begin{pmatrix}1\\-3\end{pmatrix},
  \qquad
  \ket{f_2}=\frac{1}{\sqrt{10}}\begin{pmatrix}3\\1\end{pmatrix}}
\]

两组中的每个向量范数都为 $1$，且
\[
  \braket{e_1}{e_2}=\frac{12-12}{25}=0,
  \qquad
  \braket{f_1}{f_2}=\frac{3-3}{10}=0.
\]
因此它们都是正交归一基。改变起始顺序，或将某个基向量乘以模为 $1$ 的相位因子，可以得到其他合法的正交归一基。
\end{checkedsolution}

\textbf{检查点：}Gram--Schmidt 的核心是先减去已有正交方向的投影，再对剩余向量归一化。计算中遇到复杂分数时，先提取整体公因子，归一化会将它消去。

\section{第 1.3 节教材习题与订正}
\label{exercise-set:QM-C01-S03}

\exercisemeta
  {QM-C01-S03-EX}
  {2026-08-19}
  {Shankar，第 1 章 \S 1.3，练习 1.3.1--1.3.4}
  {全部完成并逐题核对}
  {已确认（手写解答、文字作答与收口问答）}

\exercisequestion{QM-C01-S03-E01}{教材练习 1.3.1：二维 Gram--Schmidt 正交化索引}

\textbf{对应知识点：}\relatedtopic{QM-C01-S03-T04}{Gram--Schmidt 构造与完整归纳证明}。

\textbf{记录说明。}本题已在学习第 1.2 节时作为紧邻练习完成，核对后的两种构造见\relatedexercise{QM-C01-S02-E01}{教材练习 1.3.1：二维 Gram--Schmidt 正交化}。此处建立索引，不重复抄写同一解答。教材位置：印刷页 p.~15，PDF p.~29。

\exercisequestion{QM-C01-S03-E02}{教材练习 1.3.2：三维 Gram--Schmidt 正交化}

\textbf{对应知识点：}\relatedtopic{QM-C01-S03-T04}{Gram--Schmidt 构造与完整归纳证明}。

\textbf{题意概述。}从
\[
  \ket{v_1}=\begin{pmatrix}3\\0\\0\end{pmatrix},
  \qquad
  \ket{v_2}=\begin{pmatrix}0\\1\\2\end{pmatrix},
  \qquad
  \ket{v_3}=\begin{pmatrix}0\\2\\5\end{pmatrix}
\]
出发，构造教材给出的三维正交归一基。教材位置：印刷页 p.~16，PDF p.~30。

\begin{checkedsolution}
第一个单位向量为
\[
  \ket{e_1}=\frac{\ket{v_1}}{\norm{v_1}}
  =\begin{pmatrix}1\\0\\0\end{pmatrix}.
\]
由于 $\braket{e_1}{v_2}=0$，第二个剩余向量和单位向量分别为
\[
  \ket{u_2}=\ket{v_2}
  =\begin{pmatrix}0\\1\\2\end{pmatrix},
  \qquad
  \ket{e_2}=\frac{1}{\sqrt5}
  \begin{pmatrix}0\\1\\2\end{pmatrix}.
\]
对第三个向量，有
\[
  \braket{e_1}{v_3}=0,
  \qquad
  \braket{e_2}{v_3}=\frac{12}{\sqrt5}.
\]
因此
\begin{align*}
  \ket{u_3}
  &=\ket{v_3}-\ket{e_2}\braket{e_2}{v_3}\\
  &=\begin{pmatrix}0\\2\\5\end{pmatrix}
    -\begin{pmatrix}0\\12/5\\24/5\end{pmatrix}
   =\begin{pmatrix}0\\-2/5\\1/5\end{pmatrix}.
\end{align*}
归一化后
\[
  \ket{e_3}=\frac{1}{\sqrt5}
  \begin{pmatrix}0\\-2\\1\end{pmatrix}.
\]
最后，三个向量的范数均为 $1$，并且
\[
  \braket{e_1}{e_2}=\braket{e_1}{e_3}=0,
  \qquad
  \braket{e_2}{e_3}=\frac{-2+2}{5}=0.
\]
所以它们构成正交归一基。
\end{checkedsolution}

\textbf{检查点：}若原向量组线性无关，Gram--Schmidt 的剩余向量不可能为零；完整理由需要同时使用归纳假设中的正交归一性和张成空间不变性。

\exercisequestion{QM-C01-S03-E03}{教材练习 1.3.3：Schwarz 不等式的等号条件}

\textbf{对应知识点：}\relatedtopic{QM-C01-S03-T06}{Schwarz 不等式及等号条件}。

\textbf{题意概述。}判断 Schwarz 不等式何时取等号，并与实箭头的几何直觉比较。教材位置：印刷页 p.~17，PDF p.~31。

\begin{checkedsolution}
先设 $\ket{V}$ 与 $\ket{W}$ 都非零，并定义
\[
  \ket{Z}=\ket{V}
  -\frac{\braket{W}{V}}{\braket{W}{W}}\ket{W}.
\]
Schwarz 不等式的证明给出
\[
  \norm{V}^2-\frac{\abs{\braket{W}{V}}^2}{\norm{W}^2}
  =\braket{Z}{Z}\geq0.
\]
所以 Schwarz 等号成立当且仅当 $\braket{Z}{Z}=0$。由正定性，这等价于 $\ket{Z}=\ket{0}$，进而等价于
\[
  \ket{V}=\frac{\braket{W}{V}}{\braket{W}{W}}\ket{W}.
\]
也就是说，两个非零向量线性相关。反过来，若 $\ket{V}=a\ket{W}$，直接代入便得到
\[
  \abs{\braket{W}{V}}
  =\abs{a}\norm{W}^2
  =\norm{W}\norm{V}.
\]
若任一向量为零，等号也成立，且包含零向量的二元向量组线性相关。因此统一结论是：Schwarz 不等式取等号当且仅当两向量线性相关。
\end{checkedsolution}

\textbf{检查点：}Schwarz 等号只要求复线性相关，比例系数可以带任意相位；它不要求两个向量具有相同相位。

\exercisequestion{QM-C01-S03-E04}{教材练习 1.3.4：三角不等式及等号条件}

\textbf{对应知识点：}\relatedtopic{QM-C01-S03-T07}{三角不等式与零向量边界}。

\textbf{题意概述。}从 $\norm{V+W}^2$ 出发证明三角不等式，并确定取等号的条件。教材位置：印刷页 p.~17，PDF p.~31。

\begin{checkedsolution}
由内积的线性与共轭对称性，
\begin{align*}
  \norm{V+W}^2
  &=\braket{V+W}{V+W}\\
  &=\norm{V}^2+\norm{W}^2
    +\braket{V}{W}+\braket{W}{V}\\
  &=\norm{V}^2+\norm{W}^2
    +2\operatorname{Re}\braket{V}{W}\\
  &\leq\norm{V}^2+\norm{W}^2
    +2\abs{\braket{V}{W}}\\
  &\leq\norm{V}^2+\norm{W}^2
    +2\norm{V}\norm{W}\\
  &=\left(\norm{V}+\norm{W}\right)^2.
\end{align*}
两边非负，故
\[
  \norm{V+W}\leq\norm{V}+\norm{W}.
\]

若两向量均非零，等号要求上面两处不等式同时取等号。Schwarz 等号说明二者线性相关；$\operatorname{Re}z=\abs{z}$ 则要求 $z$ 为非负实数。令 $\ket{V}=a\ket{W}$，则
\[
  \braket{V}{W}=a^*\braket{W}{W},
\]
所以必须有 $a\in\real_{>0}$。因此非零情形下，等号成立当且仅当两向量同向。

固定写成 $\ket{V}=a\ket{W}$、$a\geq0$ 只能自动覆盖 $\ket{V}=\ket{0}$，不能覆盖 $\ket{W}=\ket{0}$ 而 $\ket{V}\neq\ket{0}$。包含全部零向量边界的对称表述是：存在 $\ket{U}$ 以及 $\alpha,\beta\in\real_{\geq0}$，使
\[
  \ket{V}=\alpha\ket{U},
  \qquad
  \ket{W}=\beta\ket{U}.
\]
\end{checkedsolution}

\textbf{检查点：}Schwarz 等号允许任意复比例系数；三角不等式还要求内积相位为零，因此非零向量之间的比例系数必须是正实数。
